\chapter{83}

\section{Mailand (IT), Sonntag, 2.
September}



Mantovani öffnete die Tür.

\emph{«Prego»}.\footnote{«Bitte.»}

Mit einladender Armbewegung bat er sie hinein. Er wies auf einen Stuhl
vor seinem Schreibtisch. 
Mantovani setzte seinen sanftesten Blick auf.

«Erstmal danke ich Ihnen, Signora - ?»

Sie schluckte.

«Ich kann Ihnen meinen Namen nicht verraten, weil - aber danke, dass sie sich am Sonntag Zeit nehmen, ich, ich kann wirklich nur heute - »

« - machen Sie sich keine Gedanken, unsere Einsätze richten sich nach den  Aufgaben. Die Polizei arbeitet gelegentlich auch sonntags,» Mantovani lächelte und sprach vorsichtig weiter, «es ist gut, dass sie angerufen haben. Sie haben also den Mann auf dem Video erkannt?»

Sie nickte schüchtern und sah aus, als würde sie gleich aufstehen und wieder gehen.

«Gut,» er machte eine Pause, «was möchten Sie mir erzählen?»

«Er hat nichts damit zu tun, überhaupt nichts.»

«Okay, und warum wissen Sie das?»

«Weil, also, weil er nur gemeldet hat, dass da oben auf dem Baum etwas
liegt, und weil er dachte, dass es ein Mensch sein könnte - »

Jetzt schaute sie Mantovani in die Augen. Sie hob ihre Hand und legte
sie sich über den Mund, als wolle sie sich weitere Erklärungen verbieten. 
Er nickte langsam und lehnte sich zurück.

«Sind Sie sich sicher?»

«Ja, natürlich, er, er ist ein guter Mensch. Er
weiss nicht, dass ich hier bin.»

«In welcher Verbindung stehen Sie zu ihm?»

Sie nahm ihre Hände vom Tisch und legte sie sich auf ihre Oberschenkel.

«Keine Eile, junge Frau, nehmen Sie sich Zeit.»

Sie atmete tief ein.

«Er, er ist mein Freund, wir waren zusammen da, aber bitte, das darf
keiner wissen, sonst, ich weiss nicht, das wäre nicht gut.»

Sie hatte zunehmend leiser gesprochen, umfing nun ihren Hals und
schielte nach unten auf den Tisch. Mantovani kam sie wie ein aus dem
Nest gefallenes Küken vor. Keine sechzehn, schätzte er. Es fiel ihm
leicht, einen väterlichen Ton anzuschlagen.

«Gut, seien Sie erstmal unbesorgt. Es geht hier nicht um Ihre Beziehung.
Es geht darum, dass Sie eine Meldung machen. Es wird mir helfen, zu
verstehen, was sich auf dem \emph{Roccia} in dieser Nacht abgespielt hat. Ich bin Ihnen dankbar,
wenn Sie mir alles erzählen, was Sie wissen, hm?»

Jetzt nickte sie schnell und wirkte offener.

«Wir waren in dieser Nacht unten auf der Wiese beim \emph{Roccia}. Wir
hatten was getrunken, also, ich war ein bisschen benebelt, und
plötzlich sah ich etwas Menschliches auf dem grossen Baum über uns, mein
Freund sagte dann etwas von einer Puppe oder so, - ich bin dann fast  
eingeschlafen und, und er brachte mich nach Hause. Später dann hat er
gesagt, dass er sich bei der Polizei gemeldet hat, noch in der Nacht,
aber dass das unter uns bleiben soll, weil - »

« - weil sie beide eigentlich als Paar nicht auffliegen dürfen?»

«Hm.»

Sie nickte und Tränen rollten ihr über die Wangen.
Mantovanis Strategie schien aufzugehen.

«Aber als ich, als ich das Video sah,» sie weinte nun hemmungslos, «war ich
entsetzt, ich, ich konnte nicht mehr schlafen. Und meine Freundin hat mir geraten, sie hat mir geraten, eine
Aussage zu machen, weil er, also mein Freund, nichts damit zu tun hat, es, es ist
unmöglich.»

Jetzt bekam ihre Stimme einen flehenden Unterton.

«Bitte glauben Sie mir.»

«Wissen Sie,» Mantovani lächelte, machte nachdenklich mehrere
Nickbewegungen und suchte ihren Blick, «ich denke, Ihr Freund hat die
Verletzte mit seiner Meldung vor dem Schlimmsten bewahrt.»

Überrascht wischte sie sich die Tränen weg. 
Egal, ob sie ihn als Respektsperson wahrnahm 
oder der Situation ein rasches Ende setzen wollte, Mantovani  wusste, er hatte sie soweit. 
Bereitwillig nannte sie ihm ihre Personalien.
Nach kurzem Anklopfen stand Emma Riva in der Tür. Sie begleitete die
junge Frau hinaus.

Dann sass er nachdenklich vor seinem Notizblock. Er starrte auf den Namen
und die Adresse des Mädchens. 
Wie kam sie dazu, ihn aufzusuchen? Ihm diese Geschichte zu
erzählen? War da was dran? Oder war sie womöglich geschickt worden? Und
wenn ja, von wem und warum? In welcher Beziehung stand sie wirklich zu
diesem nächtlichen Eindringling? Er gab den Namen des Mädchens bei
Google ein. Die Ergebnisse waren eindeutig. Sie lachte ihm auf einem
Bild als Zweitplatzierte eines sportlichen Wettkampfs entgegen. Ihre Schule
hatte eine Reihe von Bildern einer Siegerehrung ins Netz gestellt. Auf
einer veröffentlichten Rangliste war sie namentlich aufgeführt.
Er schaute auf die Uhr. Zeit für den Tagesaustausch. Von draussen
hörte er bereits Rossis Stimme. Der kam  mit Sönmez und 
Emma Riva plaudernd hinein. Sie setzten sich um seinen Schreibtisch.
Mantovani startete die kurze Sitzung mit Informationen zum Besuch der jungen
Frau, legte ihnen seinen Plan vor und verteilte die Aufgaben. Er grinste.

«Wenn ich mich nicht schwer täusche, kostet der Mann uns zwar unseren Sonntag, aber bis heute Abend haben wir ihn.»








